% Resumen y Abstract
% Este archivo se incluye desde main.tex con \input.
\setstretch{1.5}
\section*{Resumen}
\addcontentsline{toc}{section}{Resumen}
Los deportistas aficionados que combinan entrenamiento de fuerza con deportes de resistencia, como el running o el ciclismo, practican lo que se conoce como entrenamiento concurrente. Para progresar en ambas disciplinas sin acumular demasiada fatiga, necesitan controlar al mismo tiempo sus cargas de entrenamiento y su alimentación. Sin embargo, las aplicaciones disponibles en el mercado funcionan por separado. Una registra las pesas, otra el cardio y otra la nutrición, y ninguna cruza estos datos. Por esto, el deportista termina haciendo cálculos a mano y guiándose por sensaciones, lo que aumenta el riesgo de sobreentrenamiento y lesiones. Además, quien recién empieza no tiene un lugar centralizado ni personalizado donde informarse, y depende de contenido genérico de redes sociales.
\\\\
Este proyecto propone desarrollar un producto mínimo viable (MVP) de una aplicación móvil que reúna en un solo lugar el registro de fuerza, resistencia y nutrición, que a partir del cruce de esos datos entregue sugerencias de carga y alimentación, y que sugiera un plan de entrenamiento según el perfil de cada usuario, dejando siempre la decisión final en sus manos. La aplicación incluye además una red social deportiva, donde las sesiones, rutas y comidas que publican otros usuarios se pueden agregar con un botón a los propios módulos. Para esto se revisó la literatura sobre entrenamiento concurrente y se analizaron aplicaciones como Fitia, Kenso, Threshold y Hybrid, junto con las funciones sociales de Strava y Hevy. Luego se evaluaron tres alternativas de solución y se eligió una arquitectura por capas, que obtiene los datos de cardio desde el GPS del teléfono, desde concentradores como Apple Health y Google Health Connect o desde archivos \texttt{.fit}, y que usa inteligencia artificial para registrar comidas y entrenamientos a partir de fotos o lenguaje natural. El estudio de factibilidad mostró que el proyecto es viable en lo técnico, lo económico y lo operativo. El desarrollo sigue el modelo de prototipo evolutivo, en el que cada versión de la aplicación parte de la anterior hasta llegar al MVP.
\\\\
\textbf{Palabras clave:} entrenamiento concurrente, atleta híbrido, nutrición deportiva, aplicación móvil, cruce de datos, personalización, red social deportiva, wearables.
\newpage
\section*{Abstract}
\addcontentsline{toc}{section}{Abstract}
Amateur athletes who combine strength training with endurance sports, such as running or cycling, practice what is known as concurrent training. To progress in both disciplines without building up too much fatigue, they need to keep track of their training loads and their nutrition at the same time. However, the apps currently available work separately. One records weightlifting, another tracks cardio and another handles nutrition, and none of them cross-reference this data. As a result, athletes end up doing calculations by hand and relying on how they feel, which increases the risk of overtraining and injury. In addition, beginners have no centralized or personalized place to find information and depend on generic social media content.
\\\\
This project proposes the development of a minimum viable product (MVP) of a mobile app that brings strength, endurance and nutrition tracking together in one place and, by cross-referencing that data, offers training load and nutrition suggestions and a suggested training plan based on each user's profile, while always leaving the final decision to the user. The app also includes a sports social network, where sessions, routes and meals posted by other users can be added to one's own modules with a single tap. To this end, the literature on concurrent training was reviewed and apps such as Fitia, Kenso, Threshold and Hybrid were analyzed, along with the social features of Strava and Hevy. Three solution alternatives were then evaluated and a layered architecture was selected. It collects cardio data from the phone's GPS, from health hubs such as Apple Health and Google Health Connect, or from \texttt{.fit} files, and uses artificial intelligence to log meals and workouts from photos or natural language. The feasibility study showed that the project is technically, economically and operationally viable. Development follows the evolutionary prototyping model, in which each version of the app builds on the previous one until the MVP is reached.
\\\\
\textbf{Keywords:} concurrent training, hybrid athlete, sports nutrition, mobile application, data integration, personalization, sports social network, wearables.
\newpage
\setstretch{1.0}
